\documentclass[11pt]{beamer}
\usepackage[utf8]{inputenc}
\usepackage{amsmath, amssymb, graphicx, listings, xcolor}
\usepackage{bm}

% Presentation Theme
\usetheme{Madrid}
\usecolortheme{default}

% Code Block Styling
\lstset{
    language=Python, % Reusing Python highlighting for GMAT script readability
    basicstyle=\ttfamily\scriptsize,
    keywordstyle=\color{blue}\bfseries,
    commentstyle=\color{green!50!black},
    frame=single,
    backgroundcolor=\color{gray!10}
}

\title[Sequential Targeting in GMAT]{Astrodynamics 2026-1}
\subtitle{Advanced GMAT Architecture: \\ Sequential Targeting \& Patched Conics}
\author{Prof. Juan}
\institute{Universidad de Antioquia - Aerospace Engineering}
\date{Spring 2026}

\begin{document}

% --- Title Slide ---
\begin{frame}
    \titlepage
\end{frame}

% ==========================================
% SECTION 1: THE PROBLEM WITH ONE TARGETER
% ==========================================
\section{Mission Architecture}

% --- Slide 1 ---
\begin{frame}{The Problem with "One Giant Targeter"}
    If we try to solve the entire Earth-to-Jupiter mission inside a single Differential Corrector, we force the software to calculate how a 1-degree change in Earth parking orbit affects an arrival altitude 800 million kilometers away.
    
    \vspace{0.3cm}
    \textbf{The Consequence:}
    \begin{itemize}
        \item The Jacobian Sensitivity Matrix becomes highly non-linear and unstable.
        \item The solver requires massive iteration counts or diverges completely.
    \end{itemize}
    
    \vspace{0.3cm}
    \textbf{The Solution: Sequential Targeting}
    We mimic the \textit{Patched Conic Approximation} computationally. We use one Targeter to escape Earth, coast to the Sphere of Influence (SOI), and then use a second Targeter to navigate to Jupiter.
\end{frame}

% ==========================================
% SECTION 2: TARGET 1 - EARTH DEPARTURE
% ==========================================
\section{Target 1: Earth Escape}

% --- Slide 2 ---
\begin{frame}[fragile]{Target 1: Securing the Escape Vector}
    Our first goal is solely to leave Earth with the correct energy ($C_3$) and the correct geometric direction (Right Ascension and Declination of the Asymptote).
    
\begin{lstlisting}
% 1. Target the Escape Hyperbola
Target Target_EarthEscape {SolveMode = Solve};
    
    % Vary the True Anomaly (Injection Point)
    Vary Target_EarthEscape(Sat.TA = 180.0, {MaxStep = 10.0});
    Propagate EarthProp(Sat) {Sat.Earth.TrueAnomaly = Sat.TA};
    
    % Vary the Delta-V Magnitude
    Vary Target_EarthEscape(TJI_Burn.Element1 = 6.5, {MaxStep = 0.1});
    Maneuver TJI_Burn(Sat);
    
    % Achieve the Analytical Goals (No Jupiter propagation yet!)
    Achieve Target_EarthEscape(Sat.Earth.C3Energy = 82.5);
    Achieve Target_EarthEscape(Sat.Earth.RLA = -15.4);
    
EndTarget;
\end{lstlisting}
\end{frame}

% ==========================================
% SECTION 3: THE SOI TRANSITION
% ==========================================
\section{Crossing the Sphere of Influence}

% --- Slide 3 ---
\begin{frame}[fragile]{The Transition: Crossing the Boundary}
    Once \texttt{Target\_EarthEscape} converges, the spacecraft has the exact velocity needed. However, it is still sitting just a few hundred kilometers above Earth.
    
    \vspace{0.2cm}
    We must now propagate the spacecraft outward until it escapes Earth's gravitational grip ($\approx 925,000 \text{ km}$). We do this \textbf{outside} the Targeters using a \texttt{While} loop.
    
\begin{lstlisting}
% 2. Coast to the edge of the Sphere of Influence
% This occurs AFTER Target 1 has locked in the maneuver

While Sat.Earth.RMAG < 925000.0
    Propagate EarthProp(Sat) {Sat.ElapsedDays = 1.0};
EndWhile;

% At this exact line of code, we are officially in deep space!
\end{lstlisting}
\end{frame}

% ==========================================
% SECTION 4: TARGET 2 - MID-COURSE & ARRIVAL
% ==========================================
\section{Target 2: Deep Space \& Polar Orbit}

% --- Slide 4 ---
\begin{frame}[fragile]{Target 2: Reaching Jupiter's Poles}
    Now in the Sun-centered frame, we open our second Targeter. Here, we calculate the Deep Space Maneuvers (DSMs) needed to hit the Jupiter B-Plane and insert into a polar orbit.
    
\begin{lstlisting}
% 3. Target the Jupiter Arrival and Capture
Target Target_JupiterPolar {SolveMode = Solve};
    
    % Vary mid-course corrections (e.g., 200 days into flight)
    Propagate SunProp(Sat) {Sat.ElapsedDays = 200};
    Vary Target_JupiterPolar(DSM_Burn.Element1 = 0.5, ...);
    Vary Target_JupiterPolar(DSM_Burn.Element2 = 0.5, ...);
    Maneuver DSM_Burn(Sat);
    
    % Propagate to Jupiter and Capture
    Propagate SunProp(Sat) {Sat.Jupiter.Periapsis};
    
    % Achieve Polar Orbit conditions
    Achieve Target_JupiterPolar(Sat.Jupiter.BdotT = 0.0);
    Achieve Target_JupiterPolar(Sat.Jupiter.BdotR = 10000.0);
    
EndTarget;
\end{lstlisting}
\end{frame}

% --- Slide 5 ---
\begin{frame}{Summary of Sequential Architecture}
    \textbf{Why this makes you a better Astrodynamicist:}
    \begin{enumerate}
        \item \textbf{Decoupled Physics:} You separate the highly volatile Earth-escape physics from the deep-space cruise mechanics.
        \item \textbf{Easier Debugging:} If the script fails, you know exactly which phase broke. If Target 1 converges but Target 2 fails, you know your departure energy was fine, but your mid-course correction limits are wrong.
        \item \textbf{Real-World Accuracy:} Mission controllers do not plan the final Jupiter Orbit Insertion engine burn on the launchpad. They launch, cross the SOI, track the spacecraft, and \textit{then} calculate the mid-course corrections. Your script now behaves exactly like a real mission operations center!
    \end{enumerate}
\end{frame}

\end{document}


By teaching the script this way, you are showing your students that GMAT is not just a calculator, but a logical programming environment where they have total control over how the math is executed. Separating the `Target` blocks to mimic patched conics is exactly how professional mission designers keep their massive simulations stable!